\documentclass[11pt, a4paper, tikz,border=10pt]{report}

\usepackage[utf8]{inputenc}
\usepackage[T1]{fontenc}
\usepackage[french]{babel}

% Appel des packages dans le dossier preambule
\def\packagepath{../../../preambule}   % path du package princial

\usepackage{\packagepath/page_layout} % <--- Nouveau
\usepackage{\packagepath/config_info}
\usepackage{\packagepath/cadres_cours}
\usepackage{\packagepath/zones_reponse}
\usepackage{\packagepath/config_ds}
\usepackage{\packagepath/config_maths}

% --- PILOTAGE ---
%\def\visibleornot{visible}    % visible or invisible
\ifdefined\visibleornot\else\def\visibleornot{visible}\fi


\def\documentpath{./Sources_Latex}
\graphicspath{{./Images/}}


\def\ptsexoA{4}
\def\ptsexoB{5}
\def\ptsexoC{2}
\def\ptsexoD{4}
\def\ptsexoE{5}

\usetikzlibrary{shadows, calc}
\usepackage{dirtree}
% Définition des couleurs Firefox
\definecolor{fxframe}{RGB}{240, 240, 243}
\definecolor{fxtabactive}{RGB}{255, 255, 255}
\definecolor{fxurlbar}{RGB}{255, 255, 255}
\definecolor{fxbutton}{RGB}{180, 180, 185}
\definecolor{fxline}{RGB}{200, 200, 205}

\usepackage[edges]{forest} % Le package indispensable
\begin{document}

\def\level{Seconde}  % Classe à droite
\def\course{SNT}
\def\chaptername{\quad} % Titre au centre
\def\docname{EVAL~--~Web}        % Identifiant à gauche

\pagestyle{DS_LP}

\NomPrenom{}

\exods{\ptsexoA}

Indiquer si ces balises sont des balises d'ouverture, de fermeture ou une balise orpheline. Indiquer également ce sur quoi elles s'appliquent: (tous les termes ne sont pas à utiliser)

\fbox{lien hypertexte} \quad \fbox{tableau} \quad \fbox{liste à puces} \quad \fbox{image} \quad \fbox{titre de niveau 1} \quad \fbox{division d'une page} \quad \fbox{saut de ligne} \quad 
\fbox{liste numérotée} \quad \fbox{paragraphe} \quad \fbox{champ de saisie} \quad \fbox{titre de niveau 2} \quad \fbox{lien vers le CSS}

    \begin{tabular}{p{0.4\linewidth}cc}
    \hline
    Balise & Type & Application \\
    \hline
    \verb|<div>| & \reponseLigne[3.5cm]{\texttt{Ouverture}} & \reponseLigne[5.5cm]{division d'une page}\\
    \verb|</p>| & \reponseLigne[3.5cm]{\texttt{Fermeture}} & \reponseLigne[5.5cm]{paragraphe} \\
    \verb|<img src="image.jpg" alt="Image">| & \reponseLigne[3.5cm]{\texttt{Orpheline}} & \reponseLigne[5.5cm]{image} \\
    \verb|<a href="https://www.example.com">| & \reponseLigne[3.5cm]{\texttt{Ouverture}} & \reponseLigne[5.5cm]{lien hypertexte} \\
    \verb|<br>| & \reponseLigne[3.5cm]{\texttt{Orpheline}} & \reponseLigne[5.5cm]{saut de ligne} \\
    \verb|<h1>| & \reponseLigne[3.5cm]{\texttt{Ouverture}} & \reponseLigne[5.5cm]{titre de niveau 1} \\
    \verb|<input type="text" name="username">| & \reponseLigne[3.5cm]{\texttt{Orpheline}} & \reponseLigne[5.5cm]{champ de saisie} \\
    \verb|</ol>| & \reponseLigne[3.5cm]{\texttt{Fermeture}} & \reponseLigne[5.5cm]{liste numérotée} \\
    \hline
    \end{tabular}


    \bigskip


    \exods{\ptsexoB}

    Voici un code HTML. Corriger les erreurs présentes dans ce code:    

    \begin{lstlisting}[numbers=none]
<!DOCTYPE html>
<html>
    <title>Mon site web</title>
    <head>
        <link rel="stylesheet" href="style.css">
    </head>
    <body>
        <h1>Bienvenue sur mon site web</h1>
        <p>Ceci est un paragraphe de texte.</p>
        <img src="image.jpg" alt=Image>
        <a href="https:www.example.com">Lien vers un site</a>
        <ul>
            <li>Premier element de la liste
            <li>Second element de la liste
        </ul>
    </body>
<html>
    \end{lstlisting}

\pagebreak

 \exods{\ptsexoC}

        Dans le code de l'exercice 2, on a la ligne:

        \begin{lstlisting}[numbers=none]
<img src="image.jpg" alt="Image">
        \end{lstlisting}

        \begin{enumerate}
            \item         Que se passera-t-il si le fichier \texttt{image.jpg} n'est pas présent dans le dossier du projet? Que verra-t-on à la place de l'image?
        \begin{blocvisible}[height=1.5]
            Si le fichier \texttt{image.jpg} n'est pas présent dans le dossier du projet, le navigateur affichera le texte \verb|Image|.
        \end{blocvisible}

        \medskip

        \item Le fichier \verb|HTML| se nomme \verb|index.html| et l'image \verb|image.jpg| se trouve dans le répertoire suivant. Modifier la ligne du code HTML pour que l'image puisse être affichée correctement.
\medskip

\begin{minipage}{0.2\linewidth}
     \dirtree{%
    .1 /.
    .2 pictures.
    .3 image.jpg.
    .2 index.html.
    }
\end{minipage}\hfill{}
\begin{minipage}{0.78\linewidth}
   \begin{blocvisible}[height=2.5]
        La ligne du code HTML doit etre modifiée comme suit pour que l'image puisse être affichée correctement:

        \verb|<img src="pictures/image.jpg" alt="Image">|
   \end{blocvisible}
\end{minipage}


    \end{enumerate}

    \bigskip

    \exods{\ptsexoD}

    Voici un code CSS du fichier style.css. 

    \begin{lstlisting}[numbers=none]
body {
  background-color: lightblue; /* Couleur de fond */
  font-family: sans-serif; /* Police */
  color: red; /* Couleur du texte du body */
}

h1 {
  color: darkblue; /* Couleur du titre */
  text-align: center; /* Centrer le titre */
}
    \end{lstlisting}

    \begin{enumerate}
        \item Quel est le rôle de la balise \texttt{<link>} dans le code HTML précédent?
        
        \begin{blocvisible}[height=2]
            La balise \texttt{<link>} est utilisée pour lier un fichier CSS externe à la page HTML. Elle permet d'appliquer les styles définis dans le fichier CSS à la page web.
\end{blocvisible}
        
\medskip

        \item De quelle couleur sera le fond de la page web lorsque ce code CSS est appliqué?
        
        \begin{blocvisible}[height=1]
            Le fond de la page web sera de couleur bleu clair (lightblue) lorsque ce code CSS est appliqué.
            
        \end{blocvisible}

        \medskip

        \item De quelle couleur sera le texte du corps de la page web lorsque ce code CSS est appliqué?
        
        \begin{blocvisible}[height=1]
            Le texte du corps de la page web sera de couleur rouge (red) lorsque ce code CSS est appliqué.  
        \end{blocvisible}

        \medskip

        \item De quelle couleur seront les titres de niveau 1 (\texttt{<h1>}) de la page web lorsque ce code CSS est appliqué?
        
        \begin{blocvisible}[height=1]
            Les titres de niveau 1 (\texttt{<h1>}) de la page web seront de couleur bleu foncé (darkblue).
        \end{blocvisible}

        \medskip

        \end{enumerate}

        \bigskip

       

    \pagebreak

\exods{\ptsexoE}

Voici un code HTML, dessiner dans la zone suivante ce à quoi ressemblera la page web correspondante après l'application du code CSS \verb|style.css| de l'exercice 4:

\begin{lstlisting}[numbers=none]
<!DOCTYPE html>
<html>
    <head>
        <title>Evaluation SNT</title>
        <link rel="stylesheet" href="style.css">
    </head>
    <body>
        <h1>Bienvenue sur ma page de test</h1>
        <p>Ceci est mon evaluation.</p>
        <a href="https://www.exemple.com">Lien vers un site</a>
        <ol>
            <li> mon site web</li>
            <li> mon autre site web</li>
        </ol>
    </body>
</html>
    \end{lstlisting}


\begin{tikzpicture}
    % --- Fenêtre du navigateur Firefox ---
    \draw[rounded corners=6pt, fill=fxframe, draw=fxline, drop shadow] (0,0) rectangle (16,8.5);

    % --- Barre d'onglets ---
    \draw[fill=fxframe, draw=none] (0,7.8) rectangle (16,8.5);

    % Onglet Actif (style Firefox)
    \begin{scope}
        \clip (0.3, 7.8) -- (2.8, 7.8) -- (3.0, 8.4) -- (0.1, 8.4) -- cycle;
        \draw[fill=fxtabactive, draw=fxline, rounded corners=3pt] (0.2, 7.7) rectangle (2.9, 8.5);
    \end{scope}

    % Contenu de l'onglet (le <title>)
    \node[anchor=west, scale=0.6, font=\sffamily] at (0.8, 8.1) {\quad };
    
    % Favicon Firefox générique
    \fill[gray!50] (0.45, 8.1) circle (0.12);

    % --- Barre d'outils ---
    \draw[fill=fxframe, draw=fxline] (0, 7) -- (16, 7) -- (16, 7.8) -- (0, 7.8) -- cycle;

    % Boutons de navigation
    \draw[fxbutton, line width=1pt] (0.4, 7.4) circle (0.25); % Précédent
    \draw[fxbutton, fill=fxbutton, rounded corners=1pt] (0.4, 7.4) -- (0.5, 7.3) -- (0.3, 7.4) -- (0.5, 7.5) -- cycle;
    
    \draw[fxbutton, line width=1pt] (1.1, 7.4) circle (0.25); % Suivant
    \draw[fxbutton, fill=fxbutton, rounded corners=1pt] (1.1, 7.4) -- (1.0, 7.3) -- (1.2, 7.4) -- (1.0, 7.5) -- cycle;

    \draw[fxbutton, line width=1.2pt, rounded corners=1pt] (1.8, 7.3) -- (1.6, 7.5) -- (1.9, 7.5) -- (1.8, 7.3) -- (2.0, 7.5) -- (1.7, 7.5) -- cycle; % Actualiser (très simplifié)

    % Barre d'adresse/URL
    \draw[fill=fxurlbar, draw=fxline, rounded corners=4pt] (2.5, 7.15) rectangle (10.0, 7.65);
    \node[anchor=west, scale=0.6, font=\sffamily\color{gray}] at (2.7, 7.4) {\quad };

    % Icône cadenas
    \fill[gray!60] (2.9, 7.4) circle (0.08);

    % Bouton "Hamburger" menu
    \draw[fxbutton, line width=1.2pt] (11.4, 7.3) -- (11.7, 7.3);
    \draw[fxbutton, line width=1.2pt] (11.4, 7.4) -- (11.7, 7.4);
    \draw[fxbutton, line width=1.2pt] (11.4, 7.5) -- (11.7, 7.5);

    % --- Zone de contenu (le <body>) ---
    % C'est cette zone que tes élèves doivent remplir.
    \draw[fill=white, draw=fxline] (0, 0) rectangle (16, 7);
    
    % Annotation optionnelle pour l'exercice (à enlever si nécessaire)
    %\node[font=\sffamily\itshape, color=gray!30] at (6, 3.5) {Dessinez le contenu du <body> ici};

\end{tikzpicture}


\end{document}